% Generated by publication/build_sections.py; do not edit by hand.
\paragraph{Informal verdict.}
\textbf{A --- Complete and apparently correct}, with \textbf{high confidence}.

\paragraph{Lean coverage.}
Lean proves the explicit strictly decreasing lower bound psi\(q\)=10\textasciicircum{}\{-q\} for the Liouville number against every integer numerator and positive denominator, together with the truncation-tail estimates and a sharper antitone step bound.

\subsection{Audited informal answer}
Complete solution
\begin{enumerate}
\item Faithful restatement and two apparent typos
\end{enumerate}
Let

\[
x=\sum_{n\ge 1}10^{-n!}
   =0.110001000000000000000001\cdots .
\]

The problem asks for explicit decreasing functions

\[
\psi:\mathbb N^*\longrightarrow (0,1)
\]

such that, for every \(p\in\mathbb Z\) and every positive integer \(q\),

\[
\left\lvert x-\frac pq\right\rvert \ge \psi(q).
\]

There are two apparent typographical slips in the PDF:


In the introductory inequality for an algebraic irrational \(\alpha\), the displayed formula contains \(x\); it should contain \(\alpha\).


In both displayed inequalities, \(q\) is quantified in \(\mathbb N\), although \(p/q\) and \(\psi(q)\) require \(q\ge1\). Thus \(q\in\mathbb N^*\) is intended.


Here ``decreasing'' is understood in the standard French sense of non-increasing. A strictly decreasing example will also be given. Luc Pommeret

\begin{enumerate}
\item Truncations of the Liouville number
\end{enumerate}
For \(n\ge1\), put

\[
x_n:=\sum_{k=1}^{n}10^{-k!}
     =\frac{A_n}{Q_n},
\qquad
Q_n:=10^{n!},
\qquad
A_n:=\sum_{k=1}^{n}10^{n!-k!}\in\mathbb Z,
\]

and denote the tail by

\[
r_n:=x-x_n=\sum_{k=n+1}^{\infty}10^{-k!}.
\]

Also set

\[
T_n:=10^{-(n+1)!}.
\]

The first omitted term gives

\[
r_n>T_n.
\]

On the other hand, since the factorial exponents form a subset of the integers beginning at \((n+1)!\),

\[
r_n
<
\sum_{j=(n+1)!}^{\infty}10^{-j}
=
\frac{10}{9}\,10^{-(n+1)!}.
\]

Thus

\[
\boxed{\quad
T_n<r_n<\frac{10}{9}T_n.
\quad}
\tag{2.1}
\]

All series involved have positive terms, so these estimates require no rearrangement or limiting interchange.

\begin{enumerate}
\item The fundamental rational-separation estimate
\end{enumerate}
Let \(p\in\mathbb Z\) and \(q\in\mathbb N^*\). If \(p/q\ne x_n=A_n/Q_n\), then

\[
\left\lvert \frac pq-\frac{A_n}{Q_n}\right\rvert 
=
\frac{|pQ_n-qA_n|}{qQ_n}
\ge \frac1{qQ_n},
\]

because \(pQ_n-qA_n\) is a nonzero integer. Notice that neither \(p/q\) nor \(A_n/Q_n\) needs to be assumed reduced.
Therefore there are two cases:


if \(p/q=x_n\), then

\[
\left\lvert x-\frac pq\right\rvert =r_n;
\]



if \(p/q\ne x_n\), the reverse triangle inequality gives

\[
\left\lvert x-\frac pq\right\rvert 
\ge
\frac1{qQ_n}-r_n.
\]



Writing \(u_+:=\max(u,0)\), we consequently have, for every \(n\ge1\),

\[
\boxed{
\left\lvert x-\frac pq\right\rvert 
\ge
B_n(q):=
\min\left\{
r_n,\,
\left(\frac1{qQ_n}-r_n\right)_+
\right\}.
}
\tag{3.1}
\]

This dichotomy covers every rational \(p/q\).

\begin{enumerate}
\item An optimized explicit decreasing function
\end{enumerate}
Define

\[
\boxed{
\Psi(q):=
\sup_{n\ge1}
\min\left\{
r_n,\,
\left(\frac1{q\,10^{n!}}-r_n\right)_+
\right\}.
}
\tag{4.1}
\]

Theorem 1
The function \(\Psi:\mathbb N^*\to(0,1)\) is non-increasing and satisfies

\[
\boxed{
\left\lvert x-\frac pq\right\rvert \ge \Psi(q)
\qquad
(p\in\mathbb Z,\ q\in\mathbb N^*).
}
\tag{4.2}
\]

Proof
By (3.1), the left-hand side is at least \(B_n(q)\) for every \(n\), and hence at least their supremum.
For fixed \(n\), the map

\[
q\longmapsto
\min\left\{
r_n,\,
\left(\frac1{qQ_n}-r_n\right)_+
\right\}
\]

is non-increasing. The supremum of non-increasing functions is non-increasing, so \(\Psi\) is non-increasing.
It remains to show that \(\Psi(q)>0\). From (2.1),

\[
qQ_nr_n
<
\frac{10q}{9}\,10^{n!-(n+1)!}
=
\frac{10q}{9}\,10^{-n\,n!}
\longrightarrow0.
\]

Hence, for all sufficiently large \(n\),

\[
\frac1{qQ_n}\ge 2r_n.
\]

For such an \(n\), \(B_n(q)=r_n>0\). Thus \(\Psi(q)>0\).
Finally,

\[
\Psi(q)\le r_1<\frac1{90}<1.
\]

This proves the theorem. \(\square\)
Formula (4.1) is the strongest lower bound obtained by applying the elementary separation argument simultaneously to all decimal truncations \(x_n\).

\begin{enumerate}
\item A simpler closed-form answer
\end{enumerate}
The preceding optimized formula can be replaced by a particularly transparent step function.
For \(q\ge1\), define

\[
N(q):=
\min\left\{
n\ge1:
19q\le 9\cdot10^{n\,n!}
\right\},
\tag{5.1}
\]

and

\[
\boxed{
\psi_*(q):=10^{-(N(q)+1)!}.
}
\tag{5.2}
\]

The minimum exists because \(n\,n!\to\infty\).
Theorem 2
The function \(\psi_*\) is non-increasing and

\[
\boxed{
\left\lvert x-\frac pq\right\rvert \ge\psi_*(q)
\qquad
(p\in\mathbb Z,\ q\in\mathbb N^*).
}
\tag{5.3}
\]

Proof
Let \(n=N(q)\) and \(T_n=10^{-(n+1)!}\). Since

\[
\frac{1}{qQ_n}
=
\frac{10^{n\,n!}}q\,T_n,
\]

the definition of \(N(q)\) gives

\[
\frac{1}{qQ_n}\ge\frac{19}{9}T_n.
\]

Together with \(r_n<\frac{10}{9}T_n\), this yields

\[
\frac1{qQ_n}-r_n>T_n.
\]

We also have \(r_n>T_n\). Hence (3.1) gives

\[
\left\lvert x-\frac pq\right\rvert \ge T_n
=10^{-(N(q)+1)!}
=\psi_*(q).
\]

If \(q_1\le q_2\), then \(N(q_1)\le N(q_2)\). Therefore

\[
\psi_*(q_1)\ge\psi_*(q_2),
\]

so \(\psi_*\) is non-increasing. \(\square\)
The first ranges are

\[
\psi_*(q)=
\begin{cases}
10^{-2},&1\le q\le4,\\[2mm]
10^{-6},&5\le q\le4736,\\[2mm]
10^{-24},
&4737\le q\le
\left\lfloor\dfrac9{19}10^{18}\right\rfloor,
\end{cases}
\]

and so forth.
More generally, for every \(a\ge1\) and \(0<c\le1\),

\[
\psi_{a,c}(q)
:=
c\,10^{-a(N(q)+1)!}
\]

is another explicit admissible non-increasing function.

\begin{enumerate}
\item Strictly decreasing examples
\end{enumerate}
Because \(\Psi\) is positive and non-increasing,

\[
\boxed{
\widehat\Psi(q):=\frac{\Psi(q)}{q+1}
}
\tag{6.1}
\]

is strictly decreasing. Indeed,

\[
\widehat\Psi(q+1)
=
\frac{\Psi(q+1)}{q+2}
\le
\frac{\Psi(q)}{q+2}
<
\frac{\Psi(q)}{q+1}
=
\widehat\Psi(q).
\]

It remains an admissible lower bound because \(\widehat\Psi(q)\le\Psi(q)\).
The same construction works with \(\psi_*\):

\[
\boxed{
\widehat\psi_*(q)
:=
\frac{10^{-(N(q)+1)!}}{q+1}
}
\]

is explicit, strictly decreasing, and valid.

\begin{enumerate}
\item A particularly simple example: \(\psi(q)=10^{-q}\)
\end{enumerate}
Although much weaker than the optimized bound, the elementary function

\[
\boxed{\psi_0(q)=10^{-q}}
\tag{7.1}
\]

already answers the question and is strictly decreasing.
Proof
The case \(q=1\) is immediate. Indeed,

\[
\frac1{10}<x<\frac1{10}+\frac1{90}=\frac19<\frac12,
\]

so the closest integer is \(0\), and

\[
|x-p|>\frac1{10}=10^{-1}
\qquad(p\in\mathbb Z).
\]

Now suppose \(q\ge2\), and choose the unique \(n\ge1\) such that

\[
(n+1)!\le q<(n+2)!.
\tag{7.2}
\]

Then

\[
T_n=10^{-(n+1)!}\ge10^{-q}.
\tag{7.3}
\]

If \(n\ge2\), put

\[
a_n:=\frac{10^{n\,n!}}{(n+2)!}.
\]

We have

\[
a_2=\frac{10^4}{24}>\frac{19}{9},
\]

and \(a_n\) is strictly increasing, since

\[
\frac{a_{n+1}}{a_n}
=
\frac{
10^{(n+1)(n+1)!-n\,n!}
}{n+3}
1.
\]

Using \(q<(n+2)!\), we obtain

\[
\frac{1}{qQ_nT_n}
=
\frac{10^{n\,n!}}q
%
\frac{10^{n\,n!}}{(n+2)!}
%
\frac{19}{9}.
\]

Consequently,

\[
\frac1{qQ_n}-r_n
%
\left(\frac{19}{9}-\frac{10}{9}\right)T_n
=T_n.
\]

Both alternatives in (3.1) are therefore at least \(T_n\), and hence at least \(10^{-q}\).
It remains to check \(n=1\), corresponding to \(2\le q\le5\). Here \(Q_1=10\) and \(r_1<1/90\). For \(q=2\),

\[
\frac1{20}-r_1
%
\frac1{20}-\frac1{90}
=
\frac7{180}
%
10^{-2}.
\]

For \(3\le q\le5\),

\[
\frac1{10q}-r_1
\ge
\frac1{50}-\frac1{90}
=
\frac2{225}
%
10^{-3}
\ge10^{-q}.
\]

Since also \(r_1>10^{-2}\), the two cases in (3.1) again give the required bound. Thus

\[
\left\lvert x-\frac pq\right\rvert \ge10^{-q}
\]

for all \(p\in\mathbb Z\) and \(q\ge1\). \(\square\)

\begin{enumerate}
\item Sharpness at the canonical denominators
\end{enumerate}
The optimized function \(\Psi\) is exact along the natural sequence of truncation denominators.
Let

\[
Q_m=10^{m!},
\qquad
A_m=\sum_{k=1}^{m}10^{m!-k!}.
\]

Then

\[
\frac{A_m}{Q_m}=x_m,
\qquad
\left\lvert x-\frac{A_m}{Q_m}\right\rvert =r_m.
\tag{8.1}
\]

For \(m\ge2\),

\[
2Q_m^2r_m
<
\frac{20}{9}\,
10^{2m!-(m+1)!}
=
\frac{20}{9}\,
10^{-(m-1)m!}
<1.
\]

Hence

\[
\frac1{Q_m^2}-r_m>r_m,
\]

so the \(n=m\) term in (4.1) equals \(r_m\). Therefore

\[
\Psi(Q_m)\ge r_m.
\]

But applying (4.2) to \(p=A_m\) and \(q=Q_m\), and using (8.1), gives

\[
\Psi(Q_m)\le r_m.
\]

Thus

\[
\boxed{\Psi(10^{m!})=r_m\qquad(m\ge2).}
\tag{8.2}
\]

For the simpler function \(\psi_*\), one has \(N(Q_m)=m\). Indeed, the defining inequality holds for \(n=m\), whereas for \(n=m-1\),

\[
\frac{9\cdot10^{(m-1)(m-1)!}}{19\cdot10^{m!}}
=
\frac9{19}10^{-(m-1)!}<1.
\]

Hence

\[
\psi_*(Q_m)=T_m=10^{-(m+1)!}.
\]

Moreover,

\[
1<
\frac{r_m}{T_m}
<
1+
\frac{10}{9}
10^{-((m+2)!-(m+1)!)}
=
1+
\frac{10}{9}
10^{-(m+1)(m+1)!},
\]

so

\[
\frac{r_m}{\psi_*(Q_m)}\longrightarrow1.
\]

Thus the coefficient \(1\) in \(\psi_*\) is asymptotically optimal along the canonical denominators.
For example, at \(q=100=10^{2!}\),

\[
\frac{11}{100}=x_2,
\qquad
\left\lvert x-\frac{11}{100}\right\rvert 
=
10^{-6}+10^{-24}+10^{-120}+\cdots,
\]

while

\[
\psi_*(100)=10^{-6}.
\]


\begin{enumerate}
\item Why no fixed power-law lower bound is possible
\end{enumerate}
At \(q=Q_m=10^{m!}\), the rational \(A_m/Q_m\) satisfies

\[
\left\lvert x-\frac{A_m}{Q_m}\right\rvert 
=r_m
<
\frac{10}{9}10^{-(m+1)!}
=
\frac{10}{9}Q_m^{-(m+1)}.
\]

Given fixed \(d\ge1\) and \(C>0\), choosing \(m\) sufficiently large gives

\[
r_m<CQ_m^{-d}.
\]

Therefore no estimate

\[
\left\lvert x-\frac pq\right\rvert \ge Cq^{-d}
\]

can hold for all rational \(p/q\), exactly as expected for this Liouville number.
Hence explicit answers to Q850 include, from simplest to strongest,

\[
\boxed{\psi(q)=10^{-q}},
\]


\[
\boxed{
\psi_*(q)=10^{-(N(q)+1)!},
\quad
N(q)=\min\{n\ge1:19q\le9\cdot10^{n\,n!}\},
}
\]

and the optimized truncation envelope

\[
\boxed{
\Psi(q)=
\sup_{n\ge1}
\min\left\{
\sum_{k=n+1}^{\infty}10^{-k!},
\left(
\frac1{q10^{n!}}
-\sum_{k=n+1}^{\infty}10^{-k!}
\right)_+
\right\}.
}
\]
